% !TEX root = ../Report.tex
% ======================================================================
\chapter{Testing}
% ======================================================================

This chapter provides evidence that GiTrip was tested systematically at unit,
functional, and non-functional levels. The goal is to demonstrate correctness for
core algorithms (merge and scheduler) and reliability for user-facing workflows.

\section{Unit Testing}
GiTrip includes five unit test files with 64 test cases total, executed using
the Node.js native test framework (\texttt{node:test}) with strict assertions.

\begin{table}[H]
\centering
\caption{Summary of unit test coverage (64 test cases).}
\label{tab:unit_tests_summary}
\begin{tabular}{@{}p{0.30\linewidth}p{0.50\linewidth}p{0.12\linewidth}@{}}
\toprule
\textbf{Test file} & \textbf{Primary coverage} & \textbf{\# tests} \\
\midrule
\texttt{openhours.test.js} & Opening-hours parsing and interval validation & 11 \\
\texttt{merge.test.js}     & Three-way merge and conflict detection        & 16 \\
\texttt{planner.test.js}   & Time utilities and opening-hours scheduling   & 13 \\
\texttt{routing.test.js}   & Haversine fallback and constant-gap defaults  & 11 \\
\texttt{api-keys.test.js}  & External API connectivity (ORS, Google, GNews, Open-Meteo, Nominatim) & 13 \\
\midrule
\textbf{Total} &  & \textbf{64} \\
\bottomrule
\end{tabular}
\end{table}

\subsection{Opening Hours Parsing and Validation}
The opening-hours module tests include:
\begin{itemize}
  \item standard weekday ranges (e.g., \texttt{Mo-Fr 09:00-18:00}),
  \item wrap-around ranges (e.g., \texttt{Fr-Mo}),
  \item \texttt{off} keyword meaning closed,
  \item multi-span days with breaks,
  \item compound rules across different days,
  \item empty/null input handling.
\end{itemize}

Validation tests check whether a proposed interval is open, closed, nudged to a
next open slot, too long to fit, or exactly fits a boundary.

\subsection{Merge Conflict Detection}
Merge tests cover:
\begin{itemize}
  \item clean merges where only one side changes,
  \item deep-equality behaviour that is independent of key order,
  \item deletions (delete-vs-modify yields explicit
        \texttt{plan-stop-delete} conflict),
  \item time conflicts (\texttt{plan-stop-time}),
  \item field conflicts (\texttt{plan-stop-field}),
  \item whole-plan conflicts when similarity is low (\texttt{plan-whole}),
  \item deterministic ordering of merged days and stops.
\end{itemize}

\subsection{Planner Time Utilities and Scheduling Logic}
Planner tests cover both time-conversion functions (\texttt{toMin},
\texttt{toHHMM}, including null/empty inputs, negative clamping, and fractional
rounding) and the \texttt{applyOpeningHours} constraint engine, including:
\begin{itemize}
  \item unconstrained scheduling (no constraints; starts at proposed arrival),
  \item clamping to active period,
  \item soft desired windows (\texttt{desiredStart}, \texttt{desiredEnd}),
  \item opening-hours nudging to the next open slot,
  \item hard fixed-time windows (failure when travel arrives too late),
  \item rejection when the stop does not fit within the active window,
  \item rejection when the venue is closed all day.
\end{itemize}

\subsection{Routing Fallback Logic}
Routing tests cover the Haversine fallback estimator and constant-gap defaults
used when external routing APIs are unavailable:
\begin{itemize}
  \item identity case (same point returns zero travel time),
  \item realistic distance estimation (London to Paris within expected range),
  \item mode ordering (walking slower than cycling, cycling slower than driving),
  \item graceful handling of null or missing coordinates,
  \item constant-gap values for each transport mode.
\end{itemize}

\subsection{External API Connectivity}
API connectivity tests verify that every third-party service required by GiTrip
is correctly configured and reachable at test time:
\begin{itemize}
  \item \textbf{Presence checks} -- each required environment variable
        (\texttt{ORS\_API\_KEY}, \texttt{GOOGLE\_DIRECTIONS\_API\_KEY},
        \texttt{NEWS\_API\_KEY}) is set in \texttt{.env}.
  \item \textbf{ORS (OpenRouteService)} -- a real matrix request and a
        directions request confirm the key is accepted and returns valid
        duration data.
  \item \textbf{Google Directions} -- both driving and transit modes are
        tested; the test asserts that the response status is \texttt{OK}
        (not \texttt{REQUEST\_DENIED}), catching billing or API-enablement
        issues early.
  \item \textbf{GNews} -- a search query confirms the key returns articles
        and a travel-specific query returns relevant results.
  \item \textbf{Open-Meteo} -- both the forecast and archive endpoints are
        tested with a London coordinate, verifying that daily temperature
        and precipitation data are returned in the expected format.
  \item \textbf{Nominatim} -- a geocoding search confirms the endpoint
        returns results with latitude, longitude, and display name. A
        second test verifies that the \texttt{extratags} field (which
        carries opening-hours data from OSM) is present when requested.
\end{itemize}
These tests act as a deployment smoke test: a single failing assertion
pinpoints exactly which external service is misconfigured, rather than
requiring manual investigation of silent routing fallbacks.

\section{Functional Testing}
Functional testing verified that core workflows operate end-to-end.
Table~\ref{tab:functional_tests} summarises key cases.

\begin{table}[H]
\centering
\caption{Functional test cases for core workflows (summary).}
\label{tab:functional_tests}
\begin{tabular}{@{}p{0.52\linewidth}p{0.18\linewidth}p{0.22\linewidth}@{}}
\toprule
\textbf{Workflow} & \textbf{Result} & \textbf{Evidence} \\
\midrule
Create repo; view/edit tabs render              & Pass & Manual verification \\
Run auto-planner; commit created; redirect      & Pass & Commit history updated \\
Timeline edit; save commit; reload matches      & Pass & Snapshot persists \\
Map marker drag; recompute route; save          & Pass & Geometry updates \\
Create branch; switch branch; plan changes      & Pass & Head snapshot changes \\
Merge branches with time conflict; resolve      & Pass & Conflict UI + merge commit \\
Quick explore optimise; start repo from route   & Pass & Repo created from route \\
Export PDF view                                 & Pass & Browser print output \\
Export iCalendar (ICS)                          & Pass & Calendar import works \\
\bottomrule
\end{tabular}
\end{table}

\section{Non-functional Testing and Checks}
\subsection{Security Checks}
GiTrip includes:
\begin{itemize}
  \item Content Security Policy headers and anti-framing headers,
  \item four-tier rate limiting (global, API, geo search, and authentication
        endpoints),
  \item prepared statements for all database access,
  \item conservative input normalisation for user-provided fields,
  \item generic error messages in all user-facing responses (internal details
        logged server-side only, preventing information leakage).
\end{itemize}

These checks are not a substitute for a full penetration test, but they provide
reasonable baseline hardening for a deployed student project.

\subsection{Accessibility Checks}
Accessibility was guided by WCAG 2.2 \cite{ref_wcag22}. Checks included:
\begin{itemize}
  \item consistent labelling of controls,
  \item avoiding information conveyed only via colour,
  \item readable typography and spacing,
  \item basic keyboard navigation for non-map interactions.
\end{itemize}

A more formal audit using automated tools such as Lighthouse is recommended for
future work \cite{ref_lighthouse_overview}. The interactive map and timeline
remain difficult to fully validate for accessibility within the project scope.

\subsection{Resilience to External API Failures}
Functional tests included scenarios where routing keys were missing or external
calls failed. In these cases, GiTrip falls back to approximate travel-time
estimates and still produces a schedule (potentially with reduced accuracy). This
behaviour supports usability in constrained environments.

\section{Testing Summary}
The testing evidence demonstrates that the most error-prone components (merge
logic and scheduler constraints) are covered by automated unit tests and that
major user workflows function end-to-end. Remaining limitations are largely
usability and accessibility challenges inherent in interactive planning UIs,
rather than basic correctness failures.
